% Inbuilt themes in beamer

\documentclass[aspectratio=169]{beamer}
\usepackage[export]{adjustbox}
\usepackage{xcolor}

\definecolor{darkred}{RGB}{176, 31, 36}
\definecolor{darkgreen}{RGB}{44, 95, 47}

\usepackage{listings}
\usepackage{booktabs}
\usepackage{graphicx}

% Title page details: 
\title{Tema 4.2: Aprendizaje no supervisado} 
\subtitle{Redes genearativas adversarias}

% Theme choice:
\usetheme{Madrid}
\usecolortheme{seagull}
\useinnertheme{circles}

\author{Pablo Pérez Núñez}
 
\setbeamertemplate{navigation symbols}{}
\setbeamertemplate{author}{}

% Ajusta las dimensiones de la diapositiva

\setbeamertemplate{blocks}[rounded][shadow=false]
% \setbeamertemplate{blocks}[default]
\definecolor{myblue}{RGB}{27, 122, 194}
\definecolor{lightblue}{RGB}{55, 150, 225}
\definecolor{lightbluetwo}{RGB}{164, 209, 242}

\setbeamercolor{block title}{bg=lightblue,fg=white}
\setbeamercolor{block body}{bg=lightbluetwo,fg=black}

\setbeamercolor{title}{fg=white, bg =lightblue}
\setbeamercolor{frametitle}{fg=white, bg = lightblue}

\setbeamercolor*{palette secondary}{fg=white,bg=lightblue}

\defbeamertemplate*{footline}{infolines}
{
  \leavevmode%
  \hbox{%
  \begin{beamercolorbox}[wd=\paperwidth,ht=2.25ex,dp=1ex,left]{title in head/foot}%
    \hspace*{2ex} \usebeamerfont{title in head/foot}\insertshorttitle \ - \insertsubtitle
  \end{beamercolorbox}
  % \begin{beamercolorbox}[wd=.25\paperwidth,ht=2.25ex,dp=1ex,right]{date in head/foot}%
  %   \usebeamerfont{date in head/foot}\hspace*{2em}
  %   \insertframenumber{} / \inserttotalframenumber\hspace*{2ex} 
  %\end{beamercolorbox}
  }%
  \vskip0pt%
}

\setbeamertemplate{title page}{
  \centering
  % \usefonttheme{serif}
  \begin{minipage}{0.95\textwidth} % Ajustar el ancho según sea necesario

    \vspace{1em}
    {\color{myblue}\rule{\linewidth}{2pt}\par}
    \vspace{.5cm}
    {\bfseries \huge \textcolor{myblue}{\inserttitle}\par }
    \vspace{1em}
    {\bfseries \LARGE \textcolor{myblue}{\insertsubtitle}\par }
      
    \vspace{2.5em}

    \begin{columns}
      \column{0.6\textwidth}
      \vspace{6em}\\
      {\bfseries \large \textcolor{myblue}{Aprendizaje profundo / Deep Learning}} \\
      \vspace{1em}
      {\textcolor{myblue}{\insertauthor}}
      \column{0.4\textwidth}
      \vfill
      \centering
      \includegraphics[height=4.5em]{imgs/LogoAEPIA.pdf} \\
      \vspace{1em}
      \includegraphics[height=2.5em]{imgs/UIMP.png}
    \end{columns}

    \vspace{1em}

    {\color{myblue}\rule{\linewidth}{.5pt}\par}

  \end{minipage}
    

  %  \usefonttheme{default}
}

% Define the template for section slides
\AtBeginSection[]{
  \begin{frame}
    \vfill
    \LARGE\insertsectionhead\par
    {\color{myblue}\rule{\linewidth}{1pt}\par}
    \vfill
  \end{frame}
}
 

% MIS COSAS --------------------------------
% Bloque con margen inferior
\newenvironment{blockm}[1]{%
  \begin{block}{\textbf{#1}}%
  }{%
  \end{block}%
  \vspace{1em}%
}
% Dos imágenes
\newcommand{\twoimages}[3]{%
  \begin{columns}
    \begin{column}{0.5\textwidth}
      \centering
      \includegraphics[width=#3\textwidth]{#1}
    \end{column}
    \begin{column}{0.5\textwidth}
      \centering
      \includegraphics[width=#3\textwidth]{#2}
    \end{column}
  \end{columns}
}

% Comando para mostrar el título de la subsección al inicio de cada nueva subsección
\AtBeginSubsection[]
{
  \begin{frame}
    \vfill
    \small\insertsectionhead\par
    {\color{myblue}\rule{\linewidth}{1pt}\par}
	\LARGE\insertsubsectionhead\par
    \vfill
  \end{frame}
}

% ------------------------------------------

\begin{document}

% Title page frame
\begin{frame}[plain]
	\titlepage 
\end{frame}

\logo{}

\begin{frame}[t]{Introducción}
	\framesubtitle{Modelos generativos}
	\textbf{Aprendizaje no supervisado}:\\
	Extraer conocimiento de un conjunto de datos $x$ para el que no tenemos etiquetas $y$.\\
	\vspace{.4cm}
	\textbf{Tareas más comunes}:
	\begin{itemize}
		\item Clustering
		\item Autoencoders
		\item Modelos generativos
	\end{itemize}
	\vspace{.4cm}
	\begin{block}{\textbf{Objetivo de los modelos generativos}}
		Generar datos sintéticos pero realistas de alta dimensión $\mathbf{x} \sim p_{\theta}(\mathbf{x}),$ que se asemejen lo más posible a la distribución desconocida de datos $p(\mathbf{x})$.
	\end{block}

\end{frame}

\section{Generative Adversarial Networks (GANs)}

\begin{frame}[t]{Introducción}
	\begin{block}{\textbf{Generative Adversarial Networks (GANs)}}
		Las GANs son uno de los modelos de generación de datos de alta dimensión más populares. En ellas, \textbf{dos redes} se entrenan conjuntamente. Propuestas por Goodfellow et al. (2014).
	\end{block}
	\vspace{.5cm}
	\textbf{Están compuestas de}:
	\begin{itemize}
		\item Un \textbf{discriminador} $\textbf{D}$ que clasifica ejemplos en `reales' y `falsos'.
		\item Un \textbf{generador} $\textbf{G}$ para crear nuevos datos que `engañen' al discriminador.
	\end{itemize}
	\vspace{.5cm}
	\begin{block}{}
		Estas redes se denominan \textbf{adversarias} por que los objetivos de ambas redes son antagónicos.
	\end{block}
\end{frame}


\begin{frame}{Introducción}
	\textbf{Arquitectura:}\\
	\vspace{.3cm}
	\includegraphics[width=.75\textwidth, center]{imgs/tema4/gan/arch.pdf}\\
	\vspace{.3cm}
	\begin{block}{}
		\center \textbf{No existen dos $\textbf{D}$, el mismo discriminador evalúa ambos ejemplos.}
	\end{block}
\end{frame}

\begin{frame}{Introducción}
	\textbf{Analogía:} Falsificador de billetes ($\textbf{G}$) y Policía ($\textbf{D}$)\\
	\vspace{.3cm}
	\includegraphics[width=.75\textwidth, center]{imgs/tema4/gan/gan_ex1.pdf}\\
	\vspace{.3cm}
\end{frame}

\begin{frame}{Introducción}
	\textbf{Analogía:} Falsificador de billetes ($\textbf{G}$) y Policía ($\textbf{D}$)\\
	\vspace{.3cm}
	\includegraphics[width=.75\textwidth, center]{imgs/tema4/gan/gan_ex2.pdf}\\
	\vspace{.3cm}
\end{frame}

\begin{frame}{Introducción}
	\textbf{Analogía:} Falsificador de billetes ($\textbf{G}$) y Policía ($\textbf{D}$)\\
	\vspace{.3cm}
	\includegraphics[width=.75\textwidth, center]{imgs/tema4/gan/gan_ex3.pdf}\\
	\vspace{.3cm}
\end{frame}


\begin{frame}{Introducción}
	\textbf{Analogía:} Falsificador de billetes ($\textbf{G}$) y Policía ($\textbf{D}$)\\
	\vspace{.3cm}
	\includegraphics[width=.75\textwidth, center]{imgs/tema4/gan/gan_ex4.pdf}\\
	\vspace{.3cm}
\end{frame}



\begin{frame}[t]{Introducción}
	\framesubtitle{Notación}	
	Siendo $\mathcal{X}$ el espacio de los datos de entrada, y $d$ la dimensión del espacio latente.\\
	\vspace{.5cm}
	\textbf{Red generadora $\textbf{G}:\mathbb{R}^{d} \to \mathcal{X}$}:
	\begin{itemize}
		\item Mapea desde el espacio latente $\mathbb{R}^{d}$ al espacio de los datos de entrada $\mathcal{X}$.
		\item Entrenada para que, dado un ruido Gaussiano $\mathbf z\in\mathbb R^d$, produzca una muestra que siga la distribución de los datos de entrada.
	\end{itemize}
	\vspace{.5cm}
	\textbf{Red discriminadora $\textbf{D}:\mathcal{X} \to [0,1]$}:
	\begin{itemize}
		\item Entrenada para \textbf{clasificar} si un ejemplo proviene del conjunto real o del generador $\textbf{G}.$
	\end{itemize}
\end{frame}

\begin{frame}[t]{Introducción}
\framesubtitle{Notación}	
	Dado un conjunto de puntos `reales' $$X=\{x_1,\ldots, x_n\}$$
	si consideramos $\textbf{G}$ fijo, podemos entrenar $\textbf{D}$ generando ejemplos `falsos' $$X'=\{x'_1,\ldots, x'_n\}$$
	mediante $\mathbf x'=\textbf{G}(\mathbf z)$. Esto resultará en un conjunto de datos de clasificación binaria $$\mathfrak{D}=\{ {\underbrace{(x_1,1), \dots, (x_N,1)}_{\text{Ejemplos reales}}}, {\underbrace{(\textbf{G}(z_1),0), \dots, (\textbf{G}(z_N),0)}_{\text{Ejemplos generados}}}\} $$
	donde $\mathbf z \sim \mathcal{N} (0, 1)$.

\end{frame}

% \begin{frame}[t]{Introducción}
% \framesubtitle{Notación}	
% 	Dado un conjunto de puntos `reales' $$\mathbf{x} \sim q(\mathbf{x};\theta)$$
% 	si consideramos $\textbf{G}$ fijo, podemos entrenar $\textbf{D}$ generando ejemplos `falsos' $$\mathbf{z} \sim p(\mathbf{z})$$
% 	creando un conjunto de datos de clasificación binaria $$\mathfrak{D}=\{ {\underbrace{(x_1,1), \dots, (x_N,1)}_{\text{Ejemplos reales}}}, {\underbrace{(\textbf{G}(z_1),0), \dots, (\textbf{G}(z_N),0)}_{\text{Ejemplos generados}}}\} $$
% 	donde $q(\mathbf{x};\theta)$ es la distribución de los datos reales y $p(\mathbf{z})$ es la distribución de $\textbf{G}(Z)$ y $Z \sim \mathcal{N}(0,1)$.
% 	% Distribución de los datos reales, distribución de los datos generados por el generador G y distribución que se utiliza para generar la entrada aleatoria Z para el generador G
% \end{frame}

% \begin{frame}[t]{Entrenar discriminador}
% 	\begin{block}{}
% 		Resuelve un problema de clasificación binaria, por tanto necesita minimizar la función de pérdida \textbf{Binary Cross Entropy (BCE)}.
% 	\end{block}
% 	\vspace{.5cm}
% 	$$ \min_{\mathbf{D}} \{ - y \log \mathbf{D}(\mathbf x) - (1-y)\log(1-\mathbf{D}(\mathbf x)) \} $$
% 	% $$ \mathcal{L}(\mathbf{D})= - \frac{1}{2N} \left(\sum_{n=1}^{N} log(\mathbf{D}(x_n)) + \sum_{n=1}^{N} log(1 - \mathbf{D}(x_n))\right) $$ 
% 	Donde $\mathbf{D}(x)$ es la predicción del discriminador resultado de haber aplicado previamente una función \textbf{sigmoide} $$D(\mathbf x) = \frac{1}{(1+e^{-o})}$$ a la salida $o\in\mathbb R$ de la última capa del modelo.
% \end{frame}


% \begin{frame}[t]{Entrenar generador}
% 	Para generar cada uno de los ejemplos falsos:
% 	\begin{itemize}
% 		\item Primero obtiene $\mathbf z\in\mathbb R^d$ números aleatorios\footnote{Tradicionalmente llamamos a $\mathbf z$ variable latente y a $\mathbb R^d$ espacio latente.}, típicamente de una normal $\mathbf z \sim \mathcal{N} (0, 1)$.\\
% 		\item Aplica el generador a los valores resultantes del paso anterior $\mathbf x'=G(\mathbf z)$ .
% 	\end{itemize}
% 	\vspace{.2cm}
% 	\begin{block}{}
% 		El objetivo del generador es \textbf{engañar al discriminador} para que clasifique $\mathbf x$ como datos reales, es decir, queremos que $D( G(\mathbf z)) \approx 1$.
% 	\end{block}
% 	\vspace{.2cm}
% 	En otras palabras, para un discriminador $\textbf{D}$ dado tenemos que actualizar los pesos del generador $\textbf{G}$ para \textbf{maximizar} la cross-entropy cuando $y=0$ (ejemplo falso):
% 	$$ \max_{\mathbf{G}} \{ - \log(1-\mathbf{D}(\mathbf{G}(\mathbf z))) \} $$
% \end{frame}

% \begin{frame}[t]{Minimax Loss}
% 	Podemos transformar la maximización previa en una minimización forzando al modelo para que prediga $y=1$ en los casos generados, lo que resulta en:
% 	$$\min_{\mathbf{G}} \{ - \log(\mathbf{D}(\mathbf{G}(\mathbf z))) \}$$
% 	% Si el generador hace un trabajo perfecto, entonces $\mathbf{D}(\mathbf x')\approx 1$, por tanto la loss 
% 	% $$ \max_{\mathbf{G}} \{ - \log(1-\mathbf{D}(\mathbf{G}(\mathbf z))) \} $$
% 	% estará cercana a cero. Esto dificultará el aprendizaje del discriminador, por lo que se suele minimizar la siguiente loss: 
% 	% $$\min_{\mathbf{G}} \{ - \log(\mathbf{D}(\mathbf{G}(\mathbf z))) \}$$
% 	% que simplemente es alimentar con $\mathbf x'=G(\mathbf z)$ al discriminador pero dándole la etiqueta $y=1$.\\
% 	En resumen, $\mathbf{D}$ y $\mathbf{G}$ están jugando un juego \emph{minimax} con la función objetivo: 
% 	$$\min_{\mathbf{D}} \max_{\mathbf{G}} \{ -E_{x \sim \textrm{Data}} \log \mathbf{D}(\mathbf x) - E_{z \sim \textrm{Noise}} \log(1 - \mathbf{D}(\mathbf{G}(\mathbf z))) \}$$
% \end{frame}

\begin{frame}[t]{Entrenamiento}

	\begin{block}{}
		Lo más importante en una GAN es su proceso de entrenamiento. Cada batch o step se divide en dos pasos diferentes.
	\end{block}
	\vspace{.3cm}
	Para cada batch:
	\begin{enumerate}
		\item Optimizar el discriminador.
		\begin{itemize}
			\item Congelamos los pesos del generador.
			\item Actualizamos los pesos del discriminador para minimizar la BCE.
		\end{itemize}
		\item Optimizar el generador.
		\begin{itemize}
			\item Congelamos los pesos del discriminador.
			\item Actualizamos los pesos del generador buscando engañar al discriminador.
		\end{itemize}
	\end{enumerate}
	\vspace{.3cm}
	\begin{block}{}
		\center Este tipo de optimización se conoce como \emph{Alternating Stochastic Gradient Descent}.
	\end{block}
	
\end{frame}

\begin{frame}[t]{Entrenamiento}
	\textbf{Paso 1:} Optimizar el discriminador:\\
	\begin{enumerate}
		\item Obtenemos las muestras reales del conjunto de datos $X=\{x_1,\ldots, x_n\}$.
		\item Creamos $n$ números aleatorios $\mathbf z\in\mathbb R^d$ a partir de $\mathbf z \sim \mathcal{N} (0, 1)$.
		\item Con $\mathbf{G}$ congelado, generamos las muestras $X'=\{x'_1,\ldots, x'_n\}$ aplicando $\mathbf{G}()$ a cada $\mathbf z$.
		\item Optimizamos el discriminador mediante la Binary Cross Entropy:
		\begin{itemize}
			\item Primero obtenemos la loss para las muestras reales buscando predecir un 1. 
			$$BCE(y=1) = -\log(\mathbf{D}(\mathbf x))$$
			\item A continuación, lo mismo para las muestras generadas pero buscando predecir un 0. 
			$$BCE(y=0) = -\log(1-\mathbf{D}(\mathbf{G}(\mathbf z)))$$
			\item Optimizamos utilizando $BCE(y=1)+BCE(y=0)$.
		\end{itemize}
	\end{enumerate}
\end{frame}

\begin{frame}[t]{Entrenamiento}
	\textbf{Paso 1:} Optimizar el discriminador:\\
	\vspace{.5cm}
	\includegraphics[width=.9\textwidth, center]{imgs/tema4/gan/train_step1.pdf}\\
\end{frame}

\begin{frame}[t]{Entrenamiento}
	\textbf{Paso 2:} Optimizar el generador:\\
	\begin{enumerate}
		\item Generamos las muestras $X'=\{x'_1,\ldots, x'_n\}$ mediante $\mathbf{G}(\mathbf z)$.
		\item Con $\mathbf{D}$ congelado, optimizamos el generador mediante la Binary Cross Entropy:
		\begin{itemize}
			\item Obtenemos la loss para las muestras generadas \textbf{buscando predecir un 1}. 
			$$BCE(y=1) = -\log(\mathbf{D}(\mathbf{G}(\mathbf z)))$$
			\item Optimizamos $\mathbf{G}$ utilizando la loss calculada.
		\end{itemize}
	\end{enumerate}
	\vspace{.5cm}
	\begin{block}{Engañar al discriminador}
		Al dar ejemplos falsos y entrenar el generador para que sea capaz de generar 1, el generador aprende a engañar al discriminador.
	\end{block}

\end{frame}

\begin{frame}[t]{Entrenamiento}
	\textbf{Paso 2:} Optimizar el generador:\\
	\vspace{.5cm}
	\includegraphics[width=.9\textwidth, center]{imgs/tema4/gan/train_step2.pdf}\\
\end{frame}


% \begin{frame}{Ejemplos}
% 	\includegraphics[width=.67\textwidth, center]{imgs/tema4/gan/ex_2d_1.pdf}\\
% \end{frame}

\begin{frame}{Ejemplos}
	\includegraphics[width=.67\textwidth, center]{imgs/tema4/gan/ex_2d_2.pdf}\\
\end{frame}

% \begin{frame}{Ejemplos}
% 	\includegraphics[width=.67\textwidth, center]{imgs/tema4/gan/ex_2d_3.pdf}\\
% \end{frame}

% \begin{frame}{Ejemplos}
% 	\includegraphics[width=.67\textwidth, center]{imgs/tema4/gan/ex_2d_4.pdf}\\
% \end{frame}

\begin{frame}{Ejemplos}
	\begin{block}{Generación de imágenes}
		Figura del artículo original de Goodfellow et al., 2014 donde se presentaron las GANs. Todas son generadas salvo las de borde amarillo, que son las de entrenamiento más próximas.
	\end{block}
	\vspace{.2cm}
	\includegraphics[width=.47\textwidth, center]{imgs/tema4/gan/gan-gallery.png}\\
\end{frame}

% \begin{frame}{Ejemplos}
% 	Imágenes de habitaciones generadas por una GAN tras la primera época de entrenamiento.\\
% 	\vspace{.5cm}
% 	\includegraphics[width=.67\textwidth, center]{imgs/tema4/gan/bedrooms1.png}\\
% \end{frame}

% \begin{frame}{Ejemplos}
% 	Imágenes de habitaciones generadas por una GAN tras la quinta época de entrenamiento.\\
% 	\vspace{.5cm}
% 	\includegraphics[width=.67\textwidth, center]{imgs/tema4/gan/bedrooms2.png}\\
% \end{frame}

\begin{frame}[t]{Problemas de las GANs}
	Entrenar una red GAN tradicional puede tener los siguientes problemas:\\
	\begin{itemize}
		\item \textbf{Oscilaciones sin convergencia:} A diferencia de la minimización de loss estándar, el Alternating Stochastic Gradient Descent no tiene garantía de convergencia.
		\item \textbf{Desvanecimiento de gradientes:} Cuando el discriminador es demasiado bueno, puede que el generador no sea capaz de aprender.
		\item \textbf{Colapso de modo:} El generador puede dejar de generar diversidad en las muestras y producir siempre el mismo tipo de ejemplos.
		\item \textbf{Evaluación del rendimiento:} La evaluación del rendimiento de una GAN es difícil en la práctica. ¿Cómo de buena es una imagen generada? ¿Con quién la comparamos?.
		\item \textbf{Control de la generación:} Al generar a partir de ruido $\mathbf{z}$, ¿Cómo puedo generar ejemplos de una clase concreta?

	\end{itemize}
\end{frame}


% \subsection{GANs condicionales (cGANs)}

% \begin{frame}[t]{GANs condicionales}
% 	\begin{block}{Control de la generación}
% 		Los modelos que hemos visto hasta ahora nos permiten generar nuevas muestras a partir de un vector de ruido, pero no permiten controlar la generación.
% 	\end{block}
% 	\vspace{.3cm}
% 	Muchas aplicaciones requieren de una generación condicionada:
% 	\begin{itemize}
% 		\item Predicción del siguiente fotograma en un video.
% 		\item Rellenar partes faltantes o dañadas de una imagen.
% 		\item Transferencia del estilo artístico de una imagen a otra.
% 	\end{itemize}
% \end{frame}


% \begin{frame}[t]{GANs condicionales}
% 	\begin{block}{The Conditional GAN}
% 		Este modelo propuesto por Mirza y Osindero en 2014 añade la capacidad de generación condicionada a las GANs tradicionales. 
% 	\end{block}
% 	\vspace{.3cm}
% 	Para ello, parametriza $\mathbf{G}$ y $\mathbf{D}$ con un valor condicionante $\mathbf y$:
% 	$$\min_{\mathbf{D}} \max_{\mathbf{G}} \{ -E_{x, y \sim \textrm{Data}} \log \mathbf{D}(\mathbf x, \mathbf y) - E_{z \sim \textrm{Noise},y \sim \textrm{Data}} \log(1 - \mathbf{D}(\mathbf{G}(\mathbf z, \mathbf y), \mathbf y)) \}$$

% \end{frame}


% \begin{frame}[t]{GANs condicionales}
% 	\framesubtitle{The Conditional GAN}
% 	\textbf{Ejemplo MNIST:}\\
% 	\vspace{.3cm}
% 	Queremos ser capaces de generar imágenes de un dígito concreto:
% 	\begin{itemize}
% 		\item La $\mathbf y$ sería un vector One-Hot con la clase correspondiente al dígito.
% 		\item Tenemos que introducirlo tanto en el generador como en el discriminador.
% 	\end{itemize}
% 	\vspace{.3cm}
% 	\begin{block}{}
% 		\center \textbf{En este caso, como necesitamos datos etiquetados estamos ante un problema de aprendizaje supervisado.}
% 	\end{block}
% \end{frame}

% \begin{frame}[t]{GANs condicionales}
% 	\framesubtitle{The Conditional GAN}
% 	\textbf{Ejemplo MNIST:}\\
% 	\vspace{.5cm}
% 	\includegraphics[width=.9\textwidth, center]{imgs/tema4/gan/cGAN_g.pdf}\\
% \end{frame}

% \begin{frame}[t]{GANs condicionales}
% 	\framesubtitle{The Conditional GAN}
% 	\textbf{Ejemplo MNIST:}\\
% 	\vspace{.5cm}
% 	\includegraphics[width=.9\textwidth, center]{imgs/tema4/gan/cGAN_d.pdf}\\
% \end{frame}

% \begin{frame}[t]{GANs condicionales}
% 	\framesubtitle{The Conditional GAN}
% 	\textbf{Ejemplo MNIST:}\\
% 	\vspace{.5cm}
% 	\includegraphics[width=.9\textwidth, center]{imgs/tema4/gan/cGAN.pdf}\\
% \end{frame}

% \begin{frame}[t]{GANs condicionales}
% 	\framesubtitle{The Conditional GAN}
% 	\textbf{Ejemplo MNIST:}\\
% 	\vspace{.3cm}
% 	\includegraphics[width=.65\textwidth, center]{imgs/tema4/gan/cGAN_MNIST.png}\\
% 	\vspace{.3cm}
% 	Cada fila está condicionada a una clase y en cada columna a un vector de ruido diferente.
% \end{frame}

% \begin{frame}[t]{GANs condicionales}
% 	\begin{block}{Image to image}
% 		Existen otro tipo de problemas donde queremos generar una imagen a partir de otra dada. Isola et al. presentaron en 2016 un método de resolver estos problemas denominado \emph{pix2pix}.
% 	\end{block}
% 	\vspace{.5cm}
% 	En concreto, evaluaban su método en las siguientes tareas:
% 	\begin{itemize}
% 		\item Bordes a imágenes.
% 		\item Segmentación semántica.
% 		\item Coloreado de imágenes.
% 		\item Dia a noche.
% 	\end{itemize}

% \end{frame}

% \begin{frame}[t]{GANs condicionales}
% 	\framesubtitle{Image to image}
% 	En este caso, proponen dos arquitecturas para $\mathbf{G}$:\\
% 	\vspace{.3cm}
% 	\includegraphics[width=.5\textwidth, center]{imgs/tema4/gan/pix2pix_g.png}\\
% 	\vspace{.3cm}
% 	\begin{block}{Nota}
% 		$\mathbf{G}$ solo recibe un elemento (la imagen de origen) y la aleatoriedad proviene de un DropOut.\\
% 		En este caso denominan $y$ a la imagen generada y no $x$ como en artículos previos. 
% 	\end{block}
% \end{frame}

% \begin{frame}[t]{GANs condicionales}
% 	\framesubtitle{Image to image}
% 	El discriminador $\mathbf{D}$ es una red convolucional denominada \emph{PatchGAN}.\\
% 	\begin{block}{}
% 		En lugar de clasificar una imagen completa como real o falsa, clasifica varios parches de tamaño $NxN$ de la imagen de forma independiente y luego promedia las predicciones.
% 	\end{block}
% 	\vspace{.3cm}
% 	\includegraphics[width=.85\textwidth, center]{imgs/tema4/gan/pix2pix_patch.png}\\
% \end{frame}

% \begin{frame}[t]{GANs condicionales}
% 	\framesubtitle{Image to image}
% 	\textbf{Ejemplos:}\\
% 	\vspace{.3cm}
% 	\includegraphics[width=.85\textwidth, center]{imgs/tema4/gan/pix2pix.png}\\
% \end{frame}

% \begin{frame}[t]{GANs condicionales}
% 	\framesubtitle{Image to image}
% 	El problema principal de esta técnica es que requiere pares de ejemplos con correspondencias pixel a pixel.\\
% 	\vspace{.3cm}
% 	\begin{block}{Diferentes dominios}
% 		En ciertos casos casos tenemos ejemplos de diferentes dominios y queremos ser capaces de traducir entre ellos. Tengo una foto de manzanas y quiero cambiarlas por naranjas.
% 	\end{block}
% 	\vspace{.3cm}
% 	\begin{block}{CycleGAN}
% 		En 2017, Zhu et al. presentaron un método capaz de solventar este problema. En el proponen aprender \textbf{dos generadores}.
% 	\end{block}
% \end{frame}

% \begin{frame}[t]{GANs condicionales}
% 	\framesubtitle{Image to image}
% 	\textbf{Funcionamiento:}\\
% 	\vspace{.3cm}
% 	\includegraphics[width=.85\textwidth, center]{imgs/tema4/gan/cyclegan_explain.png}\\
% \end{frame}

% \begin{frame}[t]{GANs condicionales}
% 	\framesubtitle{Image to image}
% 	\textbf{Ejemplos:}\\
% 	\vspace{.5cm}
% 	\includegraphics[width=.9\textwidth, center]{imgs/tema4/gan/cyclegan_ex1.jpeg}\\
% \end{frame}

% \begin{frame}[t]{GANs condicionales}
% 	\framesubtitle{Image to image}
% 	\textbf{Ejemplos:}\\
% 	\vspace{.3cm}
% 	\includegraphics[width=.78\textwidth, center]{imgs/tema4/gan/cyclegan_ex2.jpeg}\\
% \end{frame}

% \subsection{Otras arquitecturas y aplicaciones}

% \begin{frame}[t]{Otras arquitecturas y aplicaciones}
% 	\framesubtitle{Introducción}
% 	\vspace{.3cm}
% 	\includegraphics[width=.45\textwidth, center]{imgs/tema4/gan/timeline.png}\\
% \end{frame}

% \begin{frame}[t]{Otras arquitecturas y aplicaciones}
% 	\textbf{Progressive growing of GANs:}\\
% 	\vspace{.7cm}
% 	\includegraphics[width=.78\textwidth, center]{imgs/tema4/gan/progressive-gan.png}\\
% \end{frame}

% \begin{frame}[t]{Otras arquitecturas y aplicaciones}
% 	\textbf{Progressive growing of GANs:}\\
% 	\vspace{.3cm}
% 	\includegraphics[width=.6\textwidth, center]{imgs/tema4/gan/progressive-gan2.png}\\
% \end{frame}

% \begin{frame}[t]{Otras arquitecturas y aplicaciones}
% 	\textbf{BigGANs:}\\
% 	\vspace{.3cm}
% 	Buscan generar imágenes de alta resolución y fidelidad mediante:
% 	\begin{itemize}
% 		\item Mecanismos de atención.
% 		\item Hinge Loss.
% 		\item Class-conditional Batch Normalization.
% 		\item Etc.
% 	\end{itemize}
% 	\vspace{.3cm}
% 	\includegraphics[width=.75\textwidth, center]{imgs/tema4/gan/biggan.png}\\
% \end{frame}

% \begin{frame}[t]{Otras arquitecturas y aplicaciones}
% 	\textbf{StyleGAN (v1):}\\
% 	\vspace{.3cm}
% 	\begin{columns}[T]
% 		\begin{column}{0.6\textwidth}
% 			\begin{itemize}
% 				\item Generador capaz de separar (automáticamente) diferentes aspectos de una imagen.
% 				\item Se basa en la idea de las Progressive Growing GANs: 
% 				\item Cada aspecto se denomina estilo: 
% 				\begin{itemize}
% 					\item Estilos más genéricos en las resoluciones más bajas.
% 					\item Detalles más finos en las resoluciones más altas
% 				\end{itemize}
% 			\end{itemize}

% 		\end{column}
% 		\begin{column}{0.4\textwidth}
% 			\includegraphics[width=.9\textwidth, center]{imgs/tema4/gan/stylegan.png}\\
% 		\end{column}
% 	\end{columns}
% \end{frame}

% \begin{frame}[t]{Otras arquitecturas y aplicaciones}
% 	\textbf{StyleGAN (v1):}\\
% 	\vspace{.5cm}
% 	\includegraphics[width=.85\textwidth, center]{imgs/tema4/gan/stylegan_inter.png}\\
% 	% \centering
% 	% \animategraphics[loop,autoplay,width=.7\textwidth]{10}{imgs/tema4/gan/styleGAN/ezgif-frame-}{001}{100}
% \end{frame}

% \begin{frame}[t]{Otras arquitecturas y aplicaciones}
% 	\textbf{Few-Shot Adversarial Learning of Realistic Neural Talking Head Models:}\\
% 	\vspace{.3cm}
% 	Origen de los \emph{DeepFakes}:
% 	\begin{itemize}
% 		\item Requiere de una imagen de referencia y otra objetivo.
% 		\item Crea una versión de la imagen referencia con la pose de la imagen objetivo.
% 		\item Ambas imágenes pueden ser de diferentes dominios.
% 	\end{itemize}
% 	\vspace{.3cm}
% 	\includegraphics[width=.9\textwidth, center]{imgs/tema4/gan/portrait_2.png}\\
% \end{frame}

% % \begin{frame}[t]{Otras arquitecturas y aplicaciones}
% % 	\textbf{Few-Shot Adversarial Learning of Realistic Neural Talking Head Models:}\\
% % 	\vspace{.3cm}
% % 	\centering
% % 	\includegraphics[width=.65\textwidth, center]{imgs/tema4/gan/portrait.png}\\
% % \end{frame}


\begin{frame}[t]{Referencias}
\begin{enumerate}
	\item \href{https://glouppe.github.io/info8010-deep-learning/?p=archives-lecture-gan.md}{\textbf{Generative Adversarial Networks, Gilles Louppe}}
	\item \href{https://fleuret.org/dlc/materials/dlc-slides-11-1-GAN.pdf}{\textbf{Generative Adversarial Networks, François Fleuret}}
	\item \href{https://d2l.ai/chapter_generative-adversarial-networks/gan.html}{\textbf{Generative Adversarial Networks, Dive into DL}}
	\item \href{https://arxiv.org/pdf/1411.1784.pdf}{\textbf{Conditional Generative Adversarial Nets}}
	\item \href{https://arxiv.org/pdf/1611.07004.pdf}{\textbf{Image-to-Image Translation with Conditional Adversarial Networks}}
	\item \href{https://arxiv.org/pdf/1703.10593.pdf}{\textbf{Unpaired Image-to-Image Translation using Cycle-Consistent Adversarial Networks}}
\end{enumerate}
\end{frame}


\end{document}
